% GENERATED FILE. Edit canonical lessons, then run npm run generate:books.
% canonical-source-hash: fnv1a64:cbb5f3ba66034b59
% canonical-lessons: ML-C313-ekkil, ML-C313-kottuva, ML-C313-lepanam, ML-C313-kasayam, ML-C313-garbhini

\chapter{Hiccup, Yawn, Ointment, Herbal decoction, Pregnant woman}
\label{ch:ml-hiccup-yawn-ointment-herbal-decoction-pregnant-woman}
\hlchaptermodality{313}

\begin{chapteropening}
\textbf{By the end of this chapter:} \emph{I can say a hiccup, a yawn, an ointment, a herbal decoction and a pregnant woman.}

Complete the last lesson of chapter 313: I can say a hiccup, a yawn, an ointment, a herbal decoction and a pregnant woman.
\end{chapteropening}

\section[\texorpdfstring{\ml{എക്കിൾ} (ekkiḷ)}{എക്കിൾ (ekkiḷ)}]{\ml{എക്കിൾ} (ekkiḷ) — a hiccup}

\label{lesson:ML-C313-ekkil}



\begin{quote}
Before the new one: say the Malayalam for an injection, then the Malayalam for an ambulance.
\end{quote}

\subsection*{You'll want to know: \ml{എക്കിൾ}}
\textbf{\ml{എക്കിൾ}} — \emph{ekkiḷ} — \textquotedblleft{}a hiccup\textquotedblright{}.

\begin{cousinweb}
\textbf{\ml{എക്കിൾ}} (\emph{ekkiḷ}) is native Dravidian, one word with Kannada \emph{ekku} and Telugu \emph{ekkili} (DEDR 772). The word imitates the sound.
\end{cousinweb}

\begin{grammarlens}[title={What you've built}]
One more word for this chapter.
\end{grammarlens}

\subsection*{Your turn}
\begin{itemize}
\raggedright
  \item \emph{Say it:} \emph{ekkiḷ}
  \item \emph{Say it:} \emph{ekkiḷ}, once more
  \item \emph{Say it:} say \emph{ekkiḷ}
  \item \emph{Recall:} say the Malayalam for an injection, then the Malayalam for an ambulance, then say \emph{ekkiḷ} again
\end{itemize}

\begin{tcolorbox}[breakable,colback=teal!4,colframe=teal!35!black,title={Before you move on}]
What does \ml{എക്കിൾ} mean? (\textquotedblleft{}A hiccup\textquotedblright{}.) Say it once more.
\end{tcolorbox}
\section[\texorpdfstring{\ml{കോട്ടുവാ} (kōṭṭuvā)}{കോട്ടുവാ (kōṭṭuvā)}]{\ml{കോട്ടുവാ} (kōṭṭuvā) — a yawn}

\label{lesson:ML-C313-kottuva}



\begin{quote}
Before the new one: say the Malayalam for an ambulance, then the Malayalam for a hiccup.
\end{quote}

\subsection*{You'll want to know: \ml{കോട്ടുവാ}}
\textbf{\ml{കോട്ടുവാ}} — \emph{kōṭṭuvā} — \textquotedblleft{}a yawn\textquotedblright{}.

\begin{cousinweb}
\textbf{\ml{കോട്ടുവാ}} (\emph{kōṭṭuvā}) is a Dravidian word. DEDR 2172 lists it as \emph{kōṭṭāvi}, \textquotedblleft{}yawn, gaping\textquotedblright{}, with Tamil \emph{koṭṭāvi}, \textquotedblleft{}yawn\textquotedblright{}, and compares them, with a query, to Kui and Kuwi words for drowsiness and dozing.
\end{cousinweb}

\begin{grammarlens}[title={What you've built}]
One more word for this chapter.
\end{grammarlens}

\subsection*{Your turn}
\begin{itemize}
\raggedright
  \item \emph{Say it:} \emph{kōṭṭuvā}
  \item \emph{Say it:} \emph{kōṭṭuvā}, once more
  \item \emph{Say it:} say \emph{kōṭṭuvā}
  \item \emph{Recall:} say the Malayalam for an ambulance, then the Malayalam for a hiccup, then say \emph{kōṭṭuvā} again
\end{itemize}

\begin{tcolorbox}[breakable,colback=teal!4,colframe=teal!35!black,title={Before you move on}]
What does \ml{കോട്ടുവാ} mean? (\textquotedblleft{}A yawn\textquotedblright{}.) Say it once more.
\end{tcolorbox}
\section[\texorpdfstring{\ml{ലേപനം} (lēpanaṁ)}{ലേപനം (lēpanaṁ)}]{\ml{ലേപനം} (lēpanaṁ) — an ointment}

\label{lesson:ML-C313-lepanam}



\begin{quote}
Before the new one: say the Malayalam for a hiccup, then the Malayalam for a yawn.
\end{quote}

\subsection*{You'll want to know: \ml{ലേപനം}}
\textbf{\ml{ലേപനം}} — \emph{lēpanaṁ} — \textquotedblleft{}an ointment\textquotedblright{}.

\begin{cousinweb}
\textbf{\ml{ലേപനം}} (\emph{lēpanaṁ}) $\leftarrow$ Sanskrit \emph{lēpana}, \textquotedblleft{}the act of smearing\textquotedblright{}, taken over whole as a tatsama, with Malayalam's noun ending \emph{-ṁ} (CDIAL 11111).
\end{cousinweb}

\begin{grammarlens}[title={What you've built}]
One more word for this chapter.
\end{grammarlens}

\subsection*{Your turn}
\begin{itemize}
\raggedright
  \item \emph{Say it:} \emph{lēpanaṁ}
  \item \emph{Say it:} \emph{lēpanaṁ}, once more
  \item \emph{Say it:} say \emph{lēpanaṁ}
  \item \emph{Recall:} say the Malayalam for a hiccup, then the Malayalam for a yawn, then say \emph{lēpanaṁ} again
\end{itemize}

\begin{tcolorbox}[breakable,colback=teal!4,colframe=teal!35!black,title={Before you move on}]
What does \ml{ലേപനം} mean? (\textquotedblleft{}An ointment\textquotedblright{}.) Say it once more.
\end{tcolorbox}
\section[\texorpdfstring{\ml{കഷായം} (kaṣāyaṁ)}{കഷായം (kaṣāyaṁ)}]{\ml{കഷായം} (kaṣāyaṁ) — a herbal decoction}

\label{lesson:ML-C313-kasayam}



\begin{quote}
Before the new one: say the Malayalam for a yawn, then the Malayalam for an ointment.
\end{quote}

\subsection*{You'll want to know: \ml{കഷായം}}
\textbf{\ml{കഷായം}} — \emph{kaṣāyaṁ} — \textquotedblleft{}a herbal decoction\textquotedblright{}.

\begin{cousinweb}
\textbf{\ml{കഷായം}} (\emph{kaṣāyaṁ}) $\leftarrow$ Sanskrit \emph{kaṣāya}, \textquotedblleft{}astringent; an astringent juice\textquotedblright{}, taken over whole as a tatsama, with Malayalam's noun ending \emph{-ṁ} (CDIAL 2974).
\end{cousinweb}

\begin{grammarlens}[title={What you've built}]
One more word for this chapter.
\end{grammarlens}

\subsection*{Your turn}
\begin{itemize}
\raggedright
  \item \emph{Say it:} \emph{kaṣāyaṁ}
  \item \emph{Say it:} \emph{kaṣāyaṁ}, once more
  \item \emph{Say it:} say \emph{kaṣāyaṁ}
  \item \emph{Recall:} say the Malayalam for a yawn, then the Malayalam for an ointment, then say \emph{kaṣāyaṁ} again
\end{itemize}

\begin{tcolorbox}[breakable,colback=teal!4,colframe=teal!35!black,title={Before you move on}]
What does \ml{കഷായം} mean? (\textquotedblleft{}A herbal decoction\textquotedblright{}.) Say it once more.
\end{tcolorbox}
\section[\texorpdfstring{\ml{ഗർഭിണി} (garbhiṇi)}{ഗർഭിണി (garbhiṇi)}]{\ml{ഗർഭിണി} (garbhiṇi) — a pregnant woman}

\label{lesson:ML-C313-garbhini}



\begin{quote}
Before the new one: say the Malayalam for an ointment, then the Malayalam for a herbal decoction.
\end{quote}

\subsection*{You'll want to know: \ml{ഗർഭിണി}}
\textbf{\ml{ഗർഭിണി}} — \emph{garbhiṇi} — \textquotedblleft{}a pregnant woman\textquotedblright{}.

\begin{cousinweb}
\textbf{\ml{ഗർഭിണി}} (\emph{garbhiṇi}) $\leftarrow$ Sanskrit \emph{garbhiṇī}, \textquotedblleft{}pregnant\textquotedblright{}, taken over whole as a tatsama (CDIAL 4062). It is built on \emph{garbha}, \textquotedblleft{}womb\textquotedblright{}.
\end{cousinweb}

\begin{grammarlens}[title={What you've built}]
One more word for this chapter.
\end{grammarlens}

\subsection*{Your turn}
\begin{itemize}
\raggedright
  \item \emph{Say it:} \emph{garbhiṇi}
  \item \emph{Say it:} \emph{garbhiṇi}, once more
  \item \emph{Say it:} say \emph{garbhiṇi}
  \item \emph{Recall:} say the Malayalam for an ointment, then the Malayalam for a herbal decoction, then say \emph{garbhiṇi} again
\end{itemize}

\begin{tcolorbox}[breakable,colback=teal!4,colframe=teal!35!black,title={Before you move on}]
What does \ml{ഗർഭിണി} mean? (\textquotedblleft{}A pregnant woman\textquotedblright{}.) Say it once more.
\end{tcolorbox}
